\documentclass[11pt]{article}
\usepackage[T1]{fontenc}
\usepackage[utf8]{inputenc}
\usepackage{palatino}
\usepackage{amsmath,amssymb,amsthm}
\usepackage{mathrsfs}
\usepackage[mathscr]{eucal}
\usepackage{geometry}
\geometry{letterpaper, margin=1in}
\usepackage{hyperref}
\usepackage{enumitem}

\newtheorem{theorem}{Theorem}[section]
\newtheorem{proposition}[theorem]{Proposition}
\newtheorem{lemma}[theorem]{Lemma}
\newtheorem{corollary}[theorem]{Corollary}
\newtheorem{definition}[theorem]{Definition}
\newtheorem{remark}[theorem]{Remark}
\newtheorem{example}[theorem]{Example}

\newcommand{\dbZ}{\mathbb{Z}}
\newcommand{\dbR}{\mathbb{R}}
\newcommand{\dbC}{\mathbb{C}}
\newcommand{\dbQ}{\mathbb{Q}}
\newcommand{\dbP}{\mathbb{P}}
\newcommand{\hb}{\mathfrak{h}}

\title{The $(3,4,\infty)$ modular family of 2-tori, completed at its three special points, is a complex structure on $S^6$}
\author{}
\date{}

\begin{document}
\maketitle

% ==================== SECTION: Front matter / Abstract-like intro ====================
\noindent\textbf{The base and the family.}
Let $\Delta = \Delta(3,4,\infty)$ be the triangle group, acting on the upper half plane $\hb$; the quotient $\hb/\Delta$ is $\dbP^1$ minus a point, and adding the cusp gives $\dbP^1$ with orbifold points of orders $3$ and $4$ and the cusp. Let $V \cong \dbZ^4$ carry automorphisms $T_1, T_2$ of orders $3$ and $4$ with $T_0 := (T_1 T_2)^{-1}$ unipotent, $(T_0 - I)^2 = 0$: this is a representation $\Delta \to \mathrm{SL}(V)$, and we let it act on the dual lattice $\Lambda = V^*$ by $A_j = (T_j^{-1})^t$ (\S2). The period functions $\tau, \mu, \beta$ on $\hb$ are an equivariant period map for it (\S3), and the quotient is a family of complex 2-tori over the punctured orbifold curve: the analogue of a universal family of abelian surfaces over a modular curve, with this rank-four lattice representation in place of the symplectic one. The very general fibre is not algebraic. The very general fibre $F$ has $\mathrm{NS}(F) = \dbZ\eta$ with $\eta$ a single indefinite class, and the algebraic dimension of the completed threefold is $a(X) = 1$ (\S9).

\medskip
\noindent\textbf{The three special fibres.}
The three conjugacy classes of maximal finite and parabolic cyclic subgroups of $\Delta$ single out the three special points, and at these three points we take the three standard completions; each is a choice, and every $\ell_0$ and every admissible $v_j$ below occurs. At an elliptic point of order $m$ the local monodromy is a torus automorphism of order $m$, and the filling is Kodaira's logarithmic transform: $(\mathrm{disc} \times \mathrm{torus})/(\dbZ/m)$, the generator acting by a rotation of the disc and by the automorphism followed by a translation; the reduced fibre is of bielliptic type (\S5). At the cusp the monodromy is unipotent of rank two, and the filling $N_0$ is Mumford's toric degeneration; its central fibre $W$ is read off the prescribed $A_2$-triangulation of $\dbR^2$, the deck action being the one in which the unimodular map $B_0 \colon \bar\Lambda \to \Lambda_{\mathrm{tor}}$ induced by $M_0 - I = (T_0^{-1})^t - I$ on $\Lambda$ enters. Modulo the lattice the triangulation has one vertex, three edges and two triangles, i.e.\ $W$ is one toric surface with hexagonal moment polygon --- the degree-six del Pezzo $\mathrm{dP}_6$ --- glued to itself along opposite sides, with two triple points (\S4). Thus there is a single family, completed in the three standard ways at the three kinds of special point of an orbifold curve with a cusp.

\medskip
\noindent\textbf{The gluing and the fundamental group.}
Each filling is glued to the global family along a collar by a fibre-preserving map, and we reglue by translations by local sections (\S6, Prop.~6.3); among these, $\pi_1$ and $H_*$ depend only on three integers $\ell_0$ (at the cusp), $\ell_1, \ell_2$ (at the elliptic points) (\S7). Let $\gamma \in V$ be the linear coordinate on $\Lambda$ whose kernel $\ker\gamma = \langle \hat u, \hat w, \hat\delta\rangle$ is the span of all $(A_j - 1)\Lambda$; then the coinvariants $\Lambda/\ker\gamma \cong \dbZ$ are detected by $\gamma$ alone, and $\ell_j = \gamma(v_j)$ for the translation vectors $v_j$. The result is $\pi_1(X) \cong \dbZ/|12\ell_0 - 4\ell_1 - 3\ell_2|$ (\S7).\footnote{\S7 proves $\pi_1(X) \cong \dbZ/|p|$, $p = 12\ell_0 - 4\ell_1 - 3\ell_2$, and computes $H_*(X;\dbZ)$ for every admissible pair $(v_1,v_2)$; the value $|p|=1$ is one point of that family.} This formula admits a reading, which the body does not use. It is the formula $|a_1 a_2 e_0 + a_2 b_1 + a_1 b_2|$ of \cite{Orl72} for the order of $H_1$ of the Seifert fibred space over $S^2(a_1,a_2)$, evaluated at $e_0 = \ell_0$ and $b_j = -\ell_j$ --- the sign coming from the convention $s_j \circ g_j = e^{-2\pi i/m_j} s_j$ for the local sections --- so here $(a_1,a_2) = (3,4)$ with Seifert invariants $(3,-\ell_1)$, $(4,-\ell_2)$ and obstruction $\ell_0$; for $(\ell_0,\ell_1,\ell_2) = (0,1,-1)$ the value is $1$ and the space in question is $S^3$ with the circle action $(z,w) \mapsto (\lambda^3 z, \lambda^4 w)$, whose orbit space is $S^2(3,4)$ \cite{Orl72}. Heuristically the twist condition says: glue so that the coinvariant circle of the fibres traces out the $(3,4)$ Seifert fibration of the three-sphere over the base.

\medskip
\noindent\textbf{The homology and the Euler number.}
The remaining torus directions contribute nothing: $\hat w$ and $\hat\delta$ are vanishing cycles at the cusp, collapsed in the hexagonal fibre, and the rest is killed by monodromy, the coinvariants of $\Lambda$ having rank one. With the twists above, $X$ has the integral homology of $S^6$ and is diffeomorphic to $S^6$ (\S7, \S8). The Euler number localises at the cusp: $e(X) = e(W) = 2$ (\S7), every other fibre being a torus or a free torus quotient and so of Euler number $0$ --- one singular fibre carries the whole Euler characteristic of $S^6$, as the twelve nodal fibres carry the $12$ of a rational elliptic surface.

\medskip
\noindent\textbf{Why the argument of \cite{CDP20} does not apply.}
That argument requires a line bundle non-torsion on every fibre and propagates a vanishing statement through each singular fibre by way of its normalisation. Here $H^2(X;\dbZ) = 0$ and $\mathrm{Pic}(X) \cong \dbC$, so every line bundle is topologically trivial; at the non-normal toric fibre the differential of the fibration itself supplies the section which the passage to the normalisation assumes away, and $R^2 f_*(T_X \otimes L) \ne 0$ for all $L$. The count of $H_1$ in \cite[Lemma 4.2]{CDP20} (after \cite[Lemma 3.2]{CDP98}) assumes trivial monodromy, whereas the monodromy representation here is non-trivial; \S10 shows that this last point is repairable for $X$, while the failure of the reduction at the non-normal fibre is decisive.

\bigskip
\noindent\textbf{Setup.}
$V := \dbZ^4$, basis $(\gamma,u,w,\delta)$; $\Lambda := V^*$, dual $(\hat\gamma,\hat u,\hat w,\hat\delta)$; $T_1, T_2 \in \mathrm{Aut}(V)$, orders $3,4$ (columns: images), and $\Pi(z)$ ($\tau,\mu,\beta$ below):
\[
T_1 = \begin{pmatrix} 1 & 0 & -6 & 2 \\ 0 & -1 & 1 & 1 \\ 0 & -1 & 0 & 1 \\ 0 & 0 & 0 & 1 \end{pmatrix}, \qquad
T_2 = \begin{pmatrix} 1 & 6 & 0 & -3 \\ 0 & 0 & -1 & 1 \\ 0 & 1 & 0 & 0 \\ 0 & 0 & 0 & 1 \end{pmatrix}, \qquad
\Pi(z) := \begin{pmatrix} 6\mu & \tau & 1 & 0 \\ \beta & \mu & 0 & 1 \end{pmatrix}.
\]
$T_0 := (T_1 T_2)^{-1} = I + N$, $N^2 = 0$, $N\gamma = Nu = 0$, $Nw = -u$, $N\delta = \gamma$; on $\Lambda$: $A_j := (T_j^{-1})^t$, $M_0 := (T_0^{-1})^t$, $\Lambda_{\mathrm{tor}} := \langle \hat w, \hat\delta\rangle$, $\bar\Lambda := \Lambda/\Lambda_{\mathrm{tor}}$, $B_0 \colon \bar\Lambda \xrightarrow{\sim} \Lambda_{\mathrm{tor}}$, $\bar{\hat\gamma} \mapsto -\hat\delta$, $\bar{\hat u} \mapsto \hat w$; $A_1, A_2$ fix $\varepsilon := \hat\gamma + 2\hat u - 4\hat w$, $\varepsilon' := \hat\gamma + 3\hat u - 3\hat w$ resp. $B := \dbP^1 \ni t$, $t_c := 1/t$, $p_1,p_2,p_0 := 0,1,\infty$, $B^\circ := B \setminus \{p_0,p_1,p_2\}$, $m_1,m_2 := 3,4$; $\pi \colon \hb_z \to B \setminus \{p_0\}$ ($\hb_z, \hb$ upper half-plane) orbifold uniformisation ($m_j$ at $p_j$, cusp $p_0$), deck group $\Delta = \langle g_1,g_2 \mid g_1^3 = g_2^4 = 1\rangle$, $g_j z_j = z_j \in \pi^{-1}(p_j)$, $g_j'(z_j) = e^{-2\pi i/m_j}$, $g_0 := (g_1 g_2)^{-1}$ parabolic; $\rho_V \colon \Delta \to \mathrm{GL}(V)$, $g_j \mapsto T_j$, $M_g := \rho_V(g)^t$; $U_r \subset \pi^{-1}\{0 < |t_c| < r\}$ ($r$ small) $\langle g_0\rangle$-invariant cusp component; holomorphic $\tau \colon \hb_z \to \hb$, $\mu,\beta \colon \hb_z \to \dbC$ ($j$ modular invariant, $j(i) = 1728$):
\[
(\tau 1)\colon j(\tau) = 1728\, t\circ\pi, \qquad
(\tau 2)\colon \tau\circ g_1 = \frac{\tau - 1}{\tau}, \ \ \tau \circ g_2 = -\frac1\tau,
\]
\[
(\tau 3)\colon \tau(z_1) = e^{i\pi/3}, \ \ \tau(z_2) = i, \qquad
(\mu 1)\colon \mu\circ g_1 = \frac{\mu}{1-\tau}, \ \ \mu\circ g_2 = \mu\left(1 + \frac{\tau}{2}\right),
\]
\[
(\mu 2)\colon \mu \text{ bounded on } U_r, \qquad
(\beta 1)\colon \beta\circ g_1 = \beta + 2 - \frac{6(1-\mu)^2}{\tau}, \ \ \beta\circ g_2 = \beta - 3 - \frac{6\mu\tau}{2},
\]
\[
(\beta 2)\colon \beta + \tau \text{ bounded on } U_r, \qquad
(\beta 3)\colon \operatorname{Im}\beta - \frac{6(\operatorname{Im}\mu)^2}{\operatorname{Im}\tau} < 0 \text{ on } \hb_z.
\]
Such $\tau,\mu,\beta$ exist by Theorem 3.4; on $\dbC^2 \ni \zeta = (\zeta_1,\zeta_2)$, $\Pi(z)$ as above with columns $\Pi(z)\hat\gamma, \Pi(z)\hat u, \Pi(z)\hat w, \Pi(z)\hat\delta$: $\lambda\cdot(z,\zeta) := (z, \zeta + \Pi(z)\lambda)$, unique $R_g(z) \in \mathrm{GL}_2(\dbC)$: $\Pi(gz) = R_g(z)\Pi(z)M_g$, $\mathcal T := (\hb_z \times \dbC^2)/\Lambda$, $\Delta$ by $\tilde g(z,\zeta) := (gz, R_g(z)\zeta)$, $J := (\mathcal T|_{\hb_z \setminus \Delta\{z_1,z_2\}})/\Delta \to B^\circ$, monodromy $A_1, A_2, M_0$ on $H_1 = \Lambda$. On $U_r$, $e^{2\pi i \tau} = t_c u_1(t_c)$, $u_1$ unit at $0$; for $s := \tau - (2\pi i)^{-1}\log u_1$ (any branch; \S4) $\Pi = [sB_0 + C(t_c) \mid I]$ in bases $(\bar{\hat\gamma},\bar{\hat u})$, $(\hat w,\hat\delta)$, $C$ holomorphic at $0$. $N' := \dbZ^3 \ni (y,y_3)$, $y \in \dbZ^2 = \Lambda_{\mathrm{tor}}$ ($e_1 \leftrightarrow \hat w$, $e_2 \leftrightarrow \hat\delta$), cocharacter lattice of $(\dbC^*)^3 \ni (x_1,x_2,t_c) \subset Y$ smooth toric, fan $\mathfrak F$ cone over $A_2$-triangulation of $\dbR^2 \times \{1\}$ (vertices $\dbZ^2 \times \{1\}$, edges $\pm e_1, \pm e_2, \pm(e_1 - e_2)$), $x_k = e^{2\pi i \zeta_k}$, $t_c$ character of $(0,0,1) \in \mathrm{Hom}(N',\dbZ)$, holomorphic on $Y$. $\Lambda$ acts freely, properly discontinuously on $Y_\epsilon := \{|t_c| < \epsilon\}$, $\epsilon$ small (4.5), by $\Psi_{\bar\lambda} := (\exp 2\pi i\, C(t_c)\bar\lambda, 1)\cdot \Phi_{\bar\lambda}$, $\Phi_{\bar\lambda}$ from $(y,y_3) \mapsto (y + y_3 B_0\bar\lambda, y_3)$; $N_0 := Y_\epsilon/\Lambda$, $W := \{t_c = 0\}/\Lambda$, $\eta \colon \mathrm{dP}_6 \to W$ normalisation, $\mathrm{dP}_6$ del Pezzo of degree $6$. For $j=1,2$: $\Delta_j \ni z_j$ $g_j$-invariant disc, coordinate $s_j$, $s_j \circ g_j = e^{-2\pi i/m_j} s_j$, $s_j^{m_j} = t_j \circ \pi$, $t_j$ coordinate at $p_j$, $N_j := (\mathcal T|_{\Delta_j})/\langle (z,\zeta) \mapsto (g_j z, R_{g_j}(z)\zeta + \Pi(g_j z) v_j/m_j)\rangle$ ($N_1$, $N_2$), $v_1 := \varepsilon$, $v_2 := -\varepsilon'$ (free; 5.4). Glue to $J$, any branch: $N_j$ by $\zeta \mapsto \zeta + (2\pi i)^{-1}(\log s_j)\Pi(z)v_j$ (\S5), $(Y_\epsilon \cap (\dbC^*)^3)/\Lambda$ on $0 < |t_c| < \epsilon$ by $(z,\zeta) \mapsto (e^{2\pi i \zeta_1}, e^{2\pi i \zeta_2}, e^{2\pi i s(z)})$. $\ell_j := \gamma(v_j)$ ($\gamma \in V = \mathrm{Hom}(\Lambda,\dbZ)$): $\ell_1 = 1$, $\ell_2 = -1$, $\ell_0 := 0$. $X := (J \sqcup N_0 \sqcup N_1 \sqcup N_2)/\!\sim$ (both gluings), $f \colon X \to B$ induced.

\begin{theorem}
$X$ is a compact connected complex $3$-manifold and $f \colon X \to B$ is a surjective holomorphic map with connected fibres.
\begin{enumerate}[label=(\arabic*)]
\item Over $B^\circ$ the map $f$ is a proper submersion whose fibres are complex $2$-tori. The fibre $f^{-1}(p_0)$ is the divisor $W \subset N_0$, which is reduced, irreducible and has normal crossings. The normalisation $\eta$ identifies opposite sides of the hexagon of $(-1)$-curves of $\mathrm{dP}_6$, the double locus is three rational curves through two triple points, and $e(W) = 2$. For $j=1,2$ one has $f^*(p_j) = m_j S_j$, where $S_j := (f^{-1}(p_j))_{\mathrm{red}}$ is a smooth bielliptic surface in $N_j$, and $\mathcal O_X(S_j)|_{S_j}$ has order $m_j$.
\item The threefold is simply connected, since $\pi_1(X) \cong \dbZ/|12\ell_0 - 4\ell_1 - 3\ell_2| = 1$, and $H_*(X;\dbZ) \cong H_*(S^6;\dbZ)$; consequently $X$ is diffeomorphic to $S^6$ (\S7ff.).
\item The algebraic dimension is $a(X) = 1$ and $\mathcal M(X) = f^*\dbC(t)$: every meromorphic function on $X$ is constant along the fibres of $f$.
\item The canonical bundle is $K_X \cong f^*\mathcal O_B(-1) \otimes \mathcal O_X(2S_2)$ and is not torsion; $c_3(X) = 2$; and $\mathrm{Aut}^0(X) \cong \dbC^*$ acts along the fibres of $f$ with fixed locus a double curve of $W$.
\end{enumerate}
\end{theorem}

\begin{remark}[on \cite{CDP20}]
The Theorem contradicts \cite[Cor.\ 2.3]{CDP20} as published. For $L \in \mathrm{Pic}(X)$ set $A := (L^* \otimes K_X)|_W$: $\eta^*A$ is trivial ($H^2(X;\dbZ)=0$, $\mathrm{Pic}(\mathrm{dP}_6)$ discrete), and for $\theta$ trivialising it $d(t_c \circ f)|_W \otimes \theta$ vanishes on the double locus, so descends to $0 \ne \sigma \in \Omega^1_X|_W \otimes A$ dying in $\widetilde\Omega^1_W \otimes A$ ($\widetilde\Omega := \Omega/\mathrm{torsion}$). By Serre--Grothendieck duality on $W$ ($\omega_W \cong K_X|_W$) and base change $R^2 f_*(T_X \otimes L) \ne 0$ for all $L$ (\S10); so hypothesis (1) of \cite[Prop.\ 2.4, p.\ 680]{CDP20} holds for no $L$. The reduction of \cite[p.\ 692]{CDP20} to $H^0(\widetilde\Omega^1_S \otimes \cdot) = 0$ on reduced fibres $S$ is stated provided $L|_S$ is non-torsion, a proviso met at $W$ for general $L$, where the conclusion $H^0(\widetilde\Omega^1_W \otimes A) = 0$ does hold; but it does not suffice. Indeed, for general $L$, $A \not\cong \mathcal O_W$ forces $H^0(W, N_W^* \otimes A) = H^0(W,A) = 0$, so the conormal term vanishes identically, while $\sigma$ gives $H^0(W, \Omega^1_W|_W \otimes A) \ne 0$: this is where the reduction fails.
\end{remark}

\newpage


\tableofcontents

\begin{center}
{\large\bfseries A compact complex threefold fibred by tori over the projective line, and the six-sphere.}
\end{center}

\subsection*{Abstract}
We construct a compact connected complex manifold $X$ of dimension three together with a surjective holomorphic map $f\colon X \to \dbP^1$ whose fibres over the complement of three points $p_0, p_1, p_2$ are complex tori of dimension two, whose fibre over $p_0$ is a reduced irreducible rational surface with normal crossings (a del Pezzo surface of degree six with the opposite sides of its anticanonical hexagon identified), and whose fibres over $p_1$ and $p_2$ are multiple, of multiplicities $3$ and $4$, with bielliptic reductions $S_1, S_2$. The family over the thrice punctured sphere is given by explicit period functions on the upper half plane attached to the $(3,4,\infty)$ orbifold; the fibre over $p_0$ is filled in by a quotient of an infinite smooth toric threefold, and the fibres over $p_1, p_2$ by logarithmic transforms. We prove that $X$ is simply connected with the integral homology of $S^6$, hence diffeomorphic to $S^6$, and we compute $a(X) = 1$ with $f$ the algebraic reduction --- the N\'eron--Severi group of the very general fibre being infinite cyclic, generated by a class $\eta$ whose associated hermitian form $H_\eta$ has signature $(1,1)$ and which is therefore represented by no positive line bundle, so that the very general fibre has algebraic dimension $0$ --- together with $f_*\mathcal O_X = \mathcal O_B$, $R^1 f_* \mathcal O_X \cong \mathcal O_B \oplus \mathcal O_B(-1)$, $R^2 f_*\mathcal O_X \cong \mathcal O_B(-1)$, $h^{0,1} = 1$, $h^{0,2} = h^{0,3} = 0$, $h^{1,0} = h^{2,0} = 0$, $K_X \cong f^*\mathcal O_B(-1) \otimes \mathcal O_X(2S_2)$, $c_1 c_2 = 0$, $c_3 = 2$, $h^0(T_X) = 1$ and $\mathrm{Aut}^0(X) \cong \dbC^*$. The existence of such an $X$ is incompatible with \cite[Cor.\ 2.3]{CDP20}; the last section locates the point at which the two accounts diverge, by showing that $R^2 f_*(T_X \otimes L) \ne 0$ for every holomorphic line bundle $L$ on $X$, a consequence of the non-normality of the normal crossings fibre $W$.

\section{Introduction.}

\subsection{The result}

Throughout, $B := \dbP^1$ carries the affine coordinate $t$, and $p_1, p_2, p_0$ are the points $t=0$, $t=1$, $t=\infty$, with $t_c := 1/t$ the coordinate at $p_0$ and $B^\circ := B \setminus \Sigma$, $\Sigma := \{p_0,p_1,p_2\}$; further $m_1 := 3$, $m_2 := 4$. We write $\mathrm{dP}_6$ for the del Pezzo surface of degree $6$, the blow up of $\dbP^2$ in three non-collinear points, whose anticanonical cycle is a hexagon of six $(-1)$-curves. A bielliptic surface is a smooth compact complex surface $(E\times F)/G$ with $E,F$ elliptic curves and $G$ a finite group of translations of $E$ acting on $F$ with $F/G \cong \dbP^1$ \cite[V.5]{BHPV}, \cite{BdF}, \cite{Ser}. We write $\mathcal M(Y)$ for the field of meromorphic functions of a compact complex space $Y$, and $a(Y) = \mathrm{tr.deg}_{\dbC}\mathcal M(Y)$ for its algebraic dimension.

\begin{theorem}[Main Theorem]
There exist a compact connected complex manifold $X$ of dimension $3$ and a surjective holomorphic map $f\colon X \to B = \dbP^1$ with connected fibres, with the following properties.
\begin{enumerate}[label=(\arabic*)]
\item $f$ is a proper holomorphic submersion over $B^\circ$, and every fibre $f^{-1}(p)$ with $p \in B^\circ$ is a complex torus of dimension $2$.
\item The fibre $W := f^{-1}(p_0)$ is reduced, irreducible, and a normal crossings divisor in $X$: at every point of $W$ the function $t_c \circ f$ is, in suitable local coordinates, $z_1$, $z_1 z_2$ or $z_1 z_2 z_3$. The normalisation of $W$ is $\mathrm{dP}_6$, and $W$ is obtained from $\mathrm{dP}_6$ by identifying the three pairs of opposite sides of its hexagon of $(-1)$-curves. The double locus $D \subset W$ consists of three smooth rational curves, each passing through both triple points of $W$ and otherwise pairwise disjoint; and $e(W) = 2$.
\item The fibres over $p_1$ and $p_2$ are multiple: as divisors, $f^*(p_1) = 3S_1$ and $f^*(p_2) = 4S_2$, where $S_j = (f^{-1}(p_j))_{\mathrm{red}}$ is a smooth bielliptic surface, and the normal bundle $\mathcal O_X(S_j)|_{S_j}$ has order exactly $m_j$ in $\mathrm{Pic}(S_j)$.
\item $X$ is simply connected and $H_*(X;\dbZ) \cong H_*(S^6;\dbZ)$; in particular $b_2(X) = b_3(X) = 0$ and $e(X) = c_3(X) = 2 = e(W)$. Consequently $X$ is diffeomorphic to the $6$-sphere.
\item $a(X) = 1$, and $f$ is the algebraic reduction of $X$: $\mathcal M(X) = f^*\mathcal M(B) = f^*\dbC(t)$, that is, every meromorphic function on $X$ is the pull back under $f$ of a rational function on $B$; moreover the very general fibre of $f$ has algebraic dimension $0$.
\item $f_*\mathcal O_X = \mathcal O_B$, $R^1 f_*\mathcal O_X \cong \mathcal O_B \oplus \mathcal O_B(-1)$ and $R^2 f_*\mathcal O_X \cong \mathcal O_B(-1)$; moreover
\[
h^{0,1}(X) = 1, \qquad h^{0,2}(X) = h^{0,3}(X) = 0, \qquad h^{1,0}(X) = h^{2,0}(X) = h^{3,0}(X) = 0.
\]
In particular $\chi(\mathcal O_X) = 0$; $\mathrm{Pic}(X) = \mathrm{Pic}^0(X) \cong H^1(X,\mathcal O_X) \cong \dbC$, since $H^1(X;\dbZ) = H^2(X;\dbZ) = 0$; and the Fr\"olicher spectral sequence of $X$ does not degenerate at $E_1$, since $b_1(X) = 0 < h^{0,1}(X)$.
\item $K_X \cong f^*\mathcal O_B(-1) \otimes \mathcal O_X(2S_2)$; in particular $K_X$ is not a torsion element of $\mathrm{Pic}(X)$, because $\mathcal O_X(4S_2) = f^*\mathcal O_B(1)$ gives $K_X^{\otimes 4k} \cong f^*\mathcal O_B(-2k)$, whose space of sections is $H^0(B,\mathcal O_B(-2k)) = 0$ for $k \ge 1$ by $f_*\mathcal O_X = \mathcal O_B$, while $h^0(\mathcal O_X) = 1$ shows that the trivial bundle does have sections. The Chern classes satisfy $c_1(X) = c_2(X) = 0$ in $H^*(X;\dbZ)$, and the Chern numbers are $c_1 c_2 = 0$ and $c_3 = 2$, whence $\chi(X,T_X) = 1$ by Hirzebruch--Riemann--Roch.
\item $h^0(X,T_X) = 1$ and $\mathrm{Aut}^0(X) \cong \dbC^*$, acting along the fibres of $f$; the generating vector field is the vertical field $\xi$ determined by the monodromy invariant vector $\hat\delta$ of the period lattice, and the fixed locus of the action is a single double curve of $W$, a smooth rational curve (cf.\ \cite[Prop.\ 2.5]{HKP}).
\end{enumerate}
\end{theorem}

\begin{corollary}
The six-sphere $S^6$ carries an integrable complex structure, namely the one obtained by transporting the complex structure of $X$ along a diffeomorphism $X \to S^6$.
\end{corollary}

\begin{remark}
Since $b_2(X) = 0$, $X$ is not K\"ahler, is not Moishezon \cite[Lemma 1.4]{CDP98}, and does not belong to Fujiki's class $\mathcal C$. The invariants in items (6)--(8) agree with the known constraints on a compact complex threefold homeomorphic to $S^6$: $h^0(T_X) = 1 \le 2$ \cite[Cor.\ 1.3]{HKP}; $h^{0,1} = h^{0,2}+1$, $H^3(X,\mathcal O_X) = 0$ and $\mathrm{Pic}(X) = \mathrm{Pic}^0(X) \cong H^1(X,\mathcal O_X)$ \cite[Cor.\ 10.2]{HKP}; $K_X$ not torsion, as forced in \cite[\S10]{HKP}; and the fixed locus of the $\dbC^*$-action of item (8), computed in \S9, is of one of the two shapes permitted by \cite[Prop.\ 2.5]{HKP}.
\end{remark}

\begin{remark}[Hodge numbers not computed here]
All the Hodge numbers of $X$ are determined by item (6) except $h^{1,1}$ and $h^{1,2} = h^{2,1}$. Since $\sum_{p,q} (-1)^{p+q} h^{p,q} = e(X)$ on any compact complex manifold (the Euler characteristic is unchanged through the Fr\"olicher spectral sequence), equivalently by the Hirzebruch--Riemann--Roch identity $\sum_p (-1)^p \chi(\Omega^p_X) = c_3(X)$ (\S9.4), the values $\chi(\mathcal O_X) = 0$, $e(X) = 2$ and the vanishings of item (6) force $h^{1,1} = h^{2,2} = h^{1,2}+1$. Only $h^{1,2}$ is left undetermined; we expect $h^{1,2} = 1$ (whence $h^{1,1} = h^{2,2} = 2$). A computation would require the direct images $R^q f_*\Omega^1_X$ and $R^q f_*\Omega^2_X$, hence a description of the sheaves $\Omega^p_{X/B}$ along the three singular fibres, where they have torsion, together with an analysis of the corresponding Leray spectral sequences in the light of the non-degeneration recorded in item (6).

The construction is explicit. Over $B^\circ$ one has a family $J \to B^\circ$ of $2$-tori whose period matrix is written down in closed form in terms of three functions $\tau, \mu, \beta$ on the upper half plane $\hb_z$ uniformising the orbifold $(B;3,4,\infty)$: $\tau$ is the elliptic modular parameter, $j(\tau) = 1728\,t$, while $\mu$ and $\beta$ solve prescribed inhomogeneous transformation laws. The local monodromies at $p_1$ and $p_2$ have orders $3$ and $4$, and the local monodromy at $p_0$ is unipotent with square-zero logarithm. The three fibres over $\Sigma$ are filled in by three local models: at $p_0$ by a quotient of an infinite smooth toric threefold built on the $A_2$-triangulation of the plane, in the spirit of Mumford's construction of degenerating abelian varieties \cite{Mum72,AMRT,KKMS}; at $p_1$ and $p_2$ by logarithmic transforms in the sense of Kodaira \cite{Kod64}, with twist vectors $v_1 = \varepsilon$ and $v_2 = -\varepsilon'$, explicit vectors of the period lattice fixed by the respective local monodromies. The three twists are what makes $X$ simply connected: \S7 gives $\pi_1(X) \cong \dbZ/|p|\dbZ$ with $p = 12\ell_0 - 4\ell_1 - 3\ell_2$, where $(\ell_0,\ell_1,\ell_2)$ records the twist data, and the above choice yields $(\ell_0,\ell_1,\ell_2) = (0,1,-1)$, so $p = -1$.
\end{remark}

\subsection{Relation to \cite{CDP20}}

The Main Theorem is incompatible with a published result. Campana, Demailly and Peternell proved in \cite{CDP98} that a compact complex threefold $X$ with $b_2(X) = 0$ and $a(X) > 0$ has $c_3(X) = e(X) = 0$, and deduced that a compact complex threefold homeomorphic to $S^6$ has algebraic dimension $0$. That proof rests on \cite[Lemma 1.5]{CDP98}, which is incorrect as stated; the corrigendum \cite{CDP20} says so and re-proves the results in two cases. \cite[Thm.\ 2.1]{CDP20} treats the case in which there exists a non-constant meromorphic map $g\colon X \dashrightarrow \dbP^1$ which is not holomorphic (this case is written out in detail in \cite[Thm.\ 3.1, Cor.\ 3.2]{LRS}, following \cite{CDP98}), while \cite[Thm.\ 2.2]{CDP20} treats the case $H^1(X,\dbZ) = H^2(X,\dbZ) = 0$, $a(X) = 1$ with holomorphic algebraic reduction $f\colon X \to C$, concluding (a) $H^i(X,T_X\otimes M) = 0$ for $i \ne 1$ and generic $M \in \mathrm{Pic}^0(X)$, (b) $\chi(X,T_X\otimes M) \le 0$ and (c) $c_3(X) \le 0$; together these give \cite[Cor.\ 2.3]{CDP20}: if $X$ is homeomorphic to $S^6$ then $a(X) = 0$. The manifold of the Main Theorem has $H^1(X,\dbZ) = H^2(X,\dbZ) = 0$, $a(X) = 1$, holomorphic algebraic reduction $f\colon X \to \dbP^1$, and $c_3(X) = 2 > 0$, so that conclusion (c) fails; so does conclusion (b), since $c_1(X) = c_2(X) = 0$ and $\mathrm{Pic}(X) = \mathrm{Pic}^0(X)$ give, by Riemann--Roch, $\chi(X,T_X\otimes M) = \tfrac12 c_3(X) = 1 > 0$ for every $M \in \mathrm{Pic}(X)$. It therefore contradicts \cite[Thm.\ 2.2]{CDP20} and \cite[Cor.\ 2.3]{CDP20}, as well as \cite[Lemma 4.2]{CDP20} (the special case $g=0$, $b_1=0$ of \cite[Lemma 3.2]{CDP98}), \cite[Cor.\ 5.2]{CDP20} and \cite[Prop.\ 5.5]{CDP20} ($=$ \cite[Thm.\ 3.1, Rem.\ 3.8(1)]{CDP98}), whose conclusions fail for $X$ (their printed proofs presuppose trivial monodromy); this failure is repairable for the purpose of the proof of Cor.\ 2.3 at $X$ (\S10), unlike the reduction at the non-normal fibre, which is decisive. It does not contradict \cite[Thm.\ 2.1]{CDP20} or \cite{LRS}, whose hypothesis is the existence of some non-constant non-holomorphic meromorphic map $X \dashrightarrow \dbP^1$, and no such map exists on our $X$: by item (5) of the Main Theorem $\mathcal M(X) = f^*\dbC(t)$, so every non-constant meromorphic map $X \dashrightarrow \dbP^1$ is of the form $h \circ f$ with $h$ a non-constant rational function on $B$, hence is holomorphic.

Section 10 is devoted to this conflict and locates the step at which the argument of \cite{CDP20} and the example part company. The key point is the fibre $W = f^{-1}(p_0)$: it is a normal crossings divisor which is not normal, the six sides of its hexagon being glued to each other in three opposite pairs (item (2) of the Main Theorem). We prove (Theorem 10.5) that for every holomorphic line bundle $L$ on $X$ one has $R^2 f_*(T_X \otimes L) \ne 0$, the obstruction being supported at $p_0$ and arising from the section $df \otimes e$ of a sheaf supported on $W$ which exists precisely because $W$ is non-normal. The vanishing of this second direct image is the hypothesis under which \cite[Prop.\ 2.4]{CDP20} is invoked in the proof of \cite[Thm.\ 2.2]{CDP20}; for $X$ it fails, for every $L$, and so the proof of \cite[Thm.\ 2.2]{CDP20} does not cover $X$. Section 10 also shows, in the light of the monodromy group of $f$ computed in \S3, that the conclusions of \cite[Lemma 4.2]{CDP20}, \cite[Cor.\ 5.2]{CDP20} and \cite[Prop.\ 5.5]{CDP20} fail for $X$ as well. Section 10 discusses the two steps of the printed proofs that do not apply to $X$; the two statements cannot both be correct. For this reason the construction is given with every matrix, chart and identification explicit.

\subsection{Plan of the paper}

Section 2 fixes the linear algebra: the lattice $V \cong \dbZ^4$ with basis $(\gamma,u,w,\delta)$, the matrices $T_1, T_2$ of orders $3$ and $4$, the unipotent $T_0 = (T_1T_2)^{-1} = I + N$ with $N^2 = 0$, and the dual picture on $\Lambda = V^*$ with $A_j = A(T_j)$ and $M_0$. All facts used later are proved there: $\Lambda_{\mathrm{tor}} = \ker(M_0 - I) = \mathrm{im}(M_0-I)$, the unimodularity of $B_0\colon \bar\Lambda \to \Lambda_{\mathrm{tor}}$, the descriptions $\Lambda^{A_1} = \langle \varepsilon,\hat\delta\rangle$, $\Lambda^{A_2} = \langle \varepsilon',\hat\delta\rangle$, and the presentation $\Delta = \langle g_1,g_2 \mid g_1^3 = g_2^4 = 1\rangle \cong \dbZ/3 * \dbZ/4$ of the orbifold fundamental group of $(B;3,4,\infty)$, the $(3,4,\infty)$ triangle group, together with the representation $\rho_V\colon \Delta \to \mathrm{GL}(V)$, $\rho_V(g_j) = T_j$. The group $\Delta$ is not isomorphic to $\mathrm{PSL}_2(\dbZ) \cong \dbZ/2 * \dbZ/3$; there is a surjection $\Delta \to \mathrm{PSL}_2(\dbZ)$, and the modular parameter $\tau$ is equivariant along it.

Section 3 constructs the period family. Theorem 3.4 establishes existence and uniqueness, up to one additive constant $c_0$ which is then chosen, of $\tau,\mu,\beta$ on $\hb_z$ subject to the equivariance and boundedness conditions listed there, their behaviour at the two elliptic points and at the cusp, the equivariance $\Pi(gz) = R_g(z)\Pi(z)M_g$ of the period matrix $\Pi = [Z \mid I]$, and the nondegeneracy $\operatorname{Im}\tau \cdot D \ne 0$ which makes $L(z) = \Pi(z)\Lambda$ a lattice in $\dbC^2$ for every $z$.\footnote{The manifold $X$ depends on the choice of $c_0$; see Remark 6.4.} This yields the family $\mathcal T \to \hb_z$ of complex $2$-tori, the quotient $J \to B^\circ$, its marking by $\Lambda$, and its monodromy.

Section 4 constructs the filling at $p_0$. The fan $\mathfrak F$ of cones over the $A_2$-triangulation of $\dbR^2 \times \{1\}$ gives a smooth toric threefold $Y_{\mathfrak F}$, not of finite type, with a global holomorphic function $t_c$; the group $\Lambda \cong \dbZ^2$ acts by the maps $\Psi_{\bar\lambda} = (c_{\bar\lambda}(t_c),1)\cdot \Phi_{\bar\lambda}$. Theorem 4.5 shows that for small radius this action is free and properly discontinuous, that $N_0$ is a complex threefold with $f_0$ proper, that the fibres over $t_c \ne 0$ are the tori of $J$ under the tautological identification $E_0$, and that $W = f_0^{-1}(0)$ is as described in item (2). For the quotient to be a complex manifold proper over the disc, only $\det B_0 \ne 0$ and the holomorphy of $C(t_c)$ at $t_c = 0$ are used; it is the unimodularity $|\det B_0| = 1$ that makes the central fibre irreducible with normalisation a single $\mathrm{dP}_6$. No positivity or symmetry of $B_0$ enters. Toric generalities are taken from \cite{Ful93,Oda88,KKMS}.

Section 5 constructs the fillings at $p_1, p_2$ by logarithmic transforms. Theorem 5.4 shows that $\tilde g_j^{\log}$, which in flat coordinates reads $(z,x) \mapsto (g_j z, A_j x + v_j/m_j)$, generates a free $\dbZ/m_j$-action on $\mathcal T|_{\Delta_j}$ --- freeness holds because $3 \nmid \gamma(v_1)$ and $\gamma(v_2)$ is odd --- that the quotient $N_j$ is a complex threefold with $f_j^*(p_j) = m_j S_j$ and $S_j$ bielliptic, and that translation by the logarithmic section $\sigma_j$ conjugates $\tilde g_j^{\log}$ to $\tilde g_j^0$, identifying $N_j \setminus f_j^{-1}(p_j)$ with $J|_{D_j^*}$.

Section 6 glues $J, N_0, N_1, N_2$ and proves Theorem 6.2: $X$ is a compact connected complex threefold, $f$ is holomorphic and surjective, and $X$ is independent, up to biholomorphism over $B$, of the auxiliary choices (radii, discs, branches of logarithm).

Section 7 computes the topology. Using the $0$-section of $J$ to split $\pi_1(J|_{B^\circ}) \cong \Lambda \rtimes \pi_1(B^\circ)$, van Kampen gives $\pi_1(X) \cong \dbZ/|p|\dbZ$ with $p = 12\ell_0 - 4\ell_1 - 3\ell_2$, hence $\pi_1(X) = 1$ here (Theorem 7.17); Mayer--Vietoris, with the homology of $W$ and of the $S_j$ computed from the cell structures of the appendix, gives $H_*(X;\dbZ) \cong H_*(S^6;\dbZ)$ and $e(X) = 2$ (Theorem 7.22). The group $H_*(X;\dbZ)$ is in fact computed twice in that section: once along the Mayer--Vietoris route just described, which uses the retraction and collapse of \S7.2, and once along the Leray route of \S7.7 together with Appendix B, which, in combination with Theorem 7.17 and Corollary 4.8, makes no use of \S7.2; the recognition of $X$ as $S^6$ therefore does not depend on Proposition 7.2. For comparison, the threefold $X'$ obtained with $v_2 = +\varepsilon'$ has $\pi_1(X') \cong \dbZ/7\dbZ$.

Section 8 concludes that $X$ is diffeomorphic to $S^6$ (Theorem 8.1): $X$ is a closed smooth simply connected $6$-manifold with the integral homology of $S^6$, hence a homotopy $6$-sphere by Hurewicz and Whitehead \cite{Hat02}, hence homeomorphic to $S^6$ by \cite{Sma61}, and diffeomorphic to $S^6$ since $\Theta_6 = 0$ \cite{KM63}, with \cite{Mil65}.

Section 9 computes the invariants of Theorem 9.1. The algebraic dimension is determined by a direct computation on the fibres: from the explicit period matrix one finds that a class $E \in \Lambda^2 V$ lies in $\mathrm{NS}(F_z)$ if and only if an explicit linear form $P_E(z)$ in the five functions $1,\tau,\mu,\beta,6\mu^2-\tau\beta$ vanishes, and these five functions are linearly independent over $\dbC$; hence for $z$ outside a countable set one gets $\mathrm{NS}(F_z) = \dbZ\eta$ with $\eta = u\wedge w + 6\,\gamma\wedge\delta$, whose hermitian form $H_\eta$ has signature $(1,1)$, so that neither $\eta$ nor $-\eta$ is the class of a positive line bundle, while $\eta^2 = 12\cdot\mathrm{vol} \ne 0$ excludes $a(F_z) = 1$ as well; the very general fibre therefore has algebraic dimension $0$. Every meromorphic function on $X$ is therefore constant on the very general fibre, whence $\mathcal M(X) = f^*\mathcal M(B) = f^*\dbC(t)$ and $a(X) = 1$, with $f$ the algebraic reduction. Section 9 also gives a second proof, by exclusion: $a(X) \ge 1$ because $f^*$ embeds $\mathcal M(B)$ into $\mathcal M(X)$; $a(X) \ne 3$ because a Moishezon threefold has $b_2 > 0$ \cite[Lemma 1.4]{CDP98} while $b_2(X) = 0$ by \S7; and $a(X) \ne 2$, since otherwise $X$ carries a meromorphic map to a curve which is not holomorphic (as in the proof of \cite[Cor.\ 2.3]{CDP20}), whence $c_3(X) = 0$ by \cite[Thm.\ 2.1]{CDP20} (see also \cite[Thm.\ 3.1]{LRS}), contradicting $c_3(X) = e(X) = 2$. This route also yields the identification $\mathcal M(X) = f^*\mathcal M(B)$, as explained in Remark 9.9: once $a(X) = 1$, the field $\mathcal M(X)$ is algebraic over the subfield $f^*\mathcal M(B)$ of transcendence degree $1$, and $f^*\mathcal M(B)$ is algebraically closed in $\mathcal M(X)$ because $f$ has connected fibres \cite[\S3]{Ue75}; hence $\mathcal M(X) = f^*\mathcal M(B)$ and $f$ is the algebraic reduction. The section then computes the direct images $R^q f_*\mathcal O_X$ and the numbers $h^{0,q}$ of item (6), the vanishing of the holomorphic forms $h^{1,0} = h^{2,0} = h^{3,0} = 0$, the canonical bundle of item (7), the Chern numbers, and the vertical $\dbC^*$-action with $h^0(T_X) = 1$ and its fixed locus. Section 10 is the discussion of \S1.2; the appendix collects the matrices, the fan $\mathfrak F$, and the cell structures used in \S7.

\subsection{Context}

The question whether $S^6$ admits a complex structure goes back to Hopf \cite{Hopf48}; for its history and the state of the art we refer to the surveys \cite{AGr,Ang18}. Among the even-dimensional spheres only $S^2$ and $S^6$ carry almost complex structures; the structure on $S^6$ induced by octonionic multiplication is not integrable, its Nijenhuis tensor being nowhere zero, and the space of almost complex structures inducing a given orientation is connected, so that no homotopy-theoretic invariant can separate a hypothetical integrable structure from the octonionic one. The one general structural theorem in this direction is that of LeBrun \cite{LeB87}, anticipated by Blanchard: no complex structure on $S^6$ is orthogonal with respect to the round metric; the proof shows that such a structure would make $S^6$ K\"ahler, contradicting $b_2 = 0$, and the later refinements also use only $H^2(S^6) = 0$. The structure constructed here is not produced from a metric and carries no a priori relation to the round one, so \cite{LeB87} does not bear on it.

Compact complex threefolds with $b_2 = 0$ are classical: Calabi and Eckmann \cite{CE53} put complex structures on $S^3\times S^3$ and on products of odd-dimensional spheres; these, the Hopf threefolds $S^1\times S^5$, the quotients $\mathrm{SL}(2,\dbC)/\Gamma$ and the principal elliptic bundles of \cite[\S4]{LRS} furnish many examples, with $a = 0,1,2$ (see \cite[\S4]{CDP98}). What distinguishes $S^6$ among these is $e = 2 \ne 0$. The complex-analytic constraints on a hypothetical complex structure on $S^6$ have been studied in detail: Huckleberry, Kebekus and Peternell \cite{HKP} proved that such an $X$ is not almost homogeneous, that $h^0(T_X) \le 2$, $H^3(X,\mathcal O_X) = 0$, $h^{0,1} = h^{0,2}+1$ and $\mathrm{Pic}(X) = \mathrm{Pic}^0(X) \cong H^1(X,\mathcal O_X)$, and analysed the possible fixed-point sets of a $\dbC^*$-action; Lehn, Rollenske and Schinko \cite{LRS} obtained bounds on spaces of sections, and wrote out in detail the meromorphic-map case of \cite{CDP98,CDP20}, following \cite{CDP98}. Constructions of complex structures on $S^6$ have been claimed by Etesi \cite{Ete15} by entirely different means, from Samelson structures on $G_2$ and conjugate orbits; nothing here depends on them or bears on their validity.

The two mechanisms combined here are classical. Logarithmic transformations were introduced by Kodaira for elliptic surfaces \cite{Kod64} (see also \cite[V.13]{BHPV}), where they are the standard device for changing the fundamental group of an elliptic fibration without changing the base; the Dolgachev surface obtained from a rational elliptic surface by logarithmic transformations of coprime multiplicities is simply connected \cite{Dol81}, and the present construction is an analogue one dimension up, with $2$-tori in place of elliptic curves. Filling in a degenerating family of abelian varieties by a toroidal model over the nilpotent monodromy cone goes back to Mumford \cite{Mum72}, and in the arithmetic setting to \cite{AMRT}; the
\newpage

\noindent combinatorial input used here is the $A_2$-triangulation of the plane. What is specific to the present situation is that the two mechanisms are used simultaneously, over a base $\dbP^1$ with exactly three critical points, for a representation $\rho_V$ of the $(3,4,\infty)$ triangle group whose local monodromies have orders $3,4$ and $\infty$, and that the resulting central fibre $W$ at the unipotent point is non-normal, with $e(W) = 2$ accounting for the whole Euler characteristic of $X$.

\section{Lattice and monodromy data.}

\subsection{The lattices $V$ and $\Lambda$}

As on page 1 (Setup), $V$ is free of rank $4$ with ordered basis $(\gamma,u,w,\delta)$ and $\Lambda := V^*$ carries the dual basis $(\hat\gamma,\hat u,\hat w,\hat\delta)$; we write $\langle \lambda, x\rangle := \lambda(x)$, elements are columns, matrices act on columns from the left, and the columns of a matrix are the images of the basis vectors. Evaluation at $\gamma$ is the functional
\[
\gamma \colon \Lambda \to \dbZ, \qquad \gamma(a\hat\gamma + b\hat u + c\hat w + d\hat\delta) = a, \qquad \ker\gamma = \langle \hat u,\hat w,\hat\delta\rangle.
\]
For $T \in \mathrm{GL}(V)$ put
\begin{equation}
\mathfrak A(T) := (T^{-1})^t \in \mathrm{GL}(\Lambda) \tag{2.1}
\end{equation}
(matrix in the dual basis): the unique element of $\mathrm{GL}(\Lambda)$ with $\langle \mathfrak A(T)\lambda, Tx\rangle = \langle \lambda,x\rangle$ for all $\lambda,x$, i.e.\ the contragredient $(T^t)^{-1}$ of $T$.

\subsection{The monodromy matrices}

\begin{definition}
In the basis $(\gamma,u,w,\delta)$ set
\[
T_1 := \begin{pmatrix} 1 & 0 & -6 & 2 \\ 0 & -1 & 1 & 1 \\ 0 & -1 & 0 & 1 \\ 0 & 0 & 0 & 1 \end{pmatrix}, \qquad
T_2 := \begin{pmatrix} 1 & 6 & 0 & -3 \\ 0 & 0 & -1 & 1 \\ 0 & 1 & 0 & 0 \\ 0 & 0 & 0 & 1 \end{pmatrix},
\]
that is,
\begin{equation}
\begin{aligned}
T_1\gamma &= \gamma, & T_1 u &= -u-w, & T_1 w &= -6\gamma+u, & T_1\delta &= 2\gamma+u+w+\delta, \\
T_2\gamma &= \gamma, & T_2 u &= 6\gamma+w, & T_2 w &= -u, & T_2\delta &= -3\gamma+u+\delta.
\end{aligned} \tag{2.2}
\end{equation}
Let $G := \langle T_1, T_2\rangle \subset \mathrm{GL}(V)$.
\end{definition}

\begin{lemma}
(i) $\det T_1 = \det T_2 = 1$. (ii) $T_1^3 = I \ne T_1$ and $T_2^4 = I \ne T_2^2$, so $T_1, T_2$ have exact orders $3,4$. (iii) Put $T_0 := (T_1T_2)^{-1}$ and $N := T_0 - I$. Then $T_0 = I+N$, $N^2 = 0$, $T_1T_2T_0 = I$, and
\[
N\gamma = Nu = 0, \qquad Nw = -u, \qquad N\delta = \gamma, \qquad \ker N = \operatorname{im} N = \langle \gamma,u\rangle;
\]
explicitly
\[
T_0 = \begin{pmatrix} 1 & 0 & 0 & 1 \\ 0 & 1 & -1 & 0 \\ 0 & 0 & 1 & 0 \\ 0 & 0 & 0 & 1 \end{pmatrix}.
\]
\end{lemma}

\begin{proof}
Direct computation from (2.2):
\[
\begin{aligned}
T_1^2 &\colon \gamma\mapsto\gamma,\ u\mapsto 6\gamma+w,\ w\mapsto -6\gamma-u-w,\ \delta\mapsto -2\gamma+u+\delta, \\
T_2^2 &\colon \gamma\mapsto\gamma,\ u\mapsto 6\gamma-u,\ w\mapsto -6\gamma-w,\ \delta\mapsto u+w+\delta, \\
T_1T_2 &\colon \gamma\mapsto\gamma,\ u\mapsto u,\ w\mapsto u+w,\ \delta\mapsto -\gamma+\delta,
\end{aligned}
\]
whence $T_1^2 \ne I \ne T_2^2$, $T_1^3 = I$, $T_2^4 = I$, and $T_1T_2 = I-N$ with $N$ as stated. Both $T_j$ fix $\gamma$, preserve $\langle \gamma,u,w\rangle$ and act trivially on $V/\langle\gamma,u,w\rangle$, and on $\langle\gamma,u,w\rangle/\langle\gamma\rangle$ by $\bar u \mapsto -\bar u - \bar w$, $\bar w \mapsto \bar u$ resp.\ $\bar u \mapsto \bar w$, $\bar w \mapsto -\bar u$, of determinant $1$; this gives (i). Since $N$ maps $w,\delta$ into $\ker N = \langle\gamma,u\rangle = \operatorname{im} N$ we get $N^2 = 0$, hence $(I-N)(I+N) = I$ and $T_0 = I+N$.
\end{proof}

\subsection{The dual matrices and the distinguished sublattices of $\Lambda$}

\begin{definition}
Put $A_j := \mathfrak A(T_j)$ ($j=1,2$) and $M_0 := \mathfrak A(T_0)$.
\end{definition}

\begin{lemma}
$\mathfrak A \colon \mathrm{GL}(V) \to \mathrm{GL}(\Lambda)$ is a group homomorphism, and in the dual basis
\[
A_1 = \begin{pmatrix} 1 & 0 & 0 & 0 \\ 6 & 0 & 1 & 0 \\ -6 & -1 & -1 & 0 \\ -2 & 1 & 0 & 1 \end{pmatrix}, \qquad
A_2 = \begin{pmatrix} 1 & 0 & 0 & 0 \\ 0 & 0 & -1 & 0 \\ -6 & 1 & 0 & 0 \\ 3 & 0 & 1 & 1 \end{pmatrix}, \qquad
M_0 = I - N^t,
\]
that is,
\begin{equation}
\begin{aligned}
A_1\hat\gamma &= \hat\gamma+6\hat u-6\hat w-2\hat\delta, & A_1\hat u &= -\hat w+\hat\delta, & A_1\hat w &= \hat u-\hat w, & A_1\hat\delta &= \hat\delta, \\
A_2\hat\gamma &= \hat\gamma-6\hat w+3\hat\delta, & A_2\hat u &= \hat w, & A_2\hat w &= -\hat u+\hat\delta, & A_2\hat\delta &= \hat\delta, \\
M_0\hat\gamma &= \hat\gamma-\hat\delta, & M_0\hat u &= \hat u+\hat w, & M_0\hat w &= \hat w, & M_0\hat\delta &= \hat\delta.
\end{aligned} \tag{2.3}
\end{equation}
Moreover $A_1A_2M_0 = I$ and $\gamma\circ A_j = \gamma \circ M_0 = \gamma$.
\end{lemma}

\begin{proof}
$\mathfrak A(ST) = ((ST)^{-1})^t = \mathfrak A(S)\mathfrak A(T)$. By Lemma 2.2(ii), $A_1 = (T_1^2)^t$ and $A_2 = (T_2^3)^t$, and $T_2^3\colon \gamma\mapsto\gamma$, $u\mapsto -w$, $w\mapsto -6\gamma+u$, $\delta\mapsto 3\gamma+w+\delta$; transposing gives the displayed $A_1, A_2$. From $T_0^{-1} = I-N$ we get $M_0 = I-N^t$, and $N^t\hat\gamma = \hat\delta$, $N^t\hat u = -\hat w$, $N^t\hat w = N^t\hat\delta = 0$ by Lemma 2.2(iii), which is (2.3). Then $A_1A_2M_0 = \mathfrak A(T_1T_2T_0) = I$, and $\gamma(A_j\lambda) = \langle \mathfrak A(T_j)\lambda, T_j\gamma\rangle = \gamma(\lambda)$ since $T_j\gamma = \gamma$; likewise for $M_0$.
\end{proof}

\begin{definition}
Put $\Lambda_{\mathrm{tor}} := \langle \hat w,\hat\delta\rangle \subset \Lambda$ and $\bar\Lambda := \Lambda/\Lambda_{\mathrm{tor}}$ with basis $(\bar{\hat\gamma},\bar{\hat u})$; identify $\bar\Lambda \cong \dbZ^2$ by $a\bar{\hat\gamma}+b\bar{\hat u} \mapsto (a,b)^t$ and $\Lambda_{\mathrm{tor}} \cong \dbZ^2$ by $c\hat w+d\hat\delta \mapsto (c,d)^t$. Set $\varepsilon := \hat\gamma+2\hat u-4\hat w$ and $\varepsilon' := \hat\gamma+3\hat u-3\hat w$.
\end{definition}

\begin{lemma}\label{lem:B0}
$\Lambda_{\mathrm{tor}} = \ker(M_0-I) = \operatorname{im}(M_0-I)$, and $M_0-I$ induces
\[
B_0 \colon \bar\Lambda \to \Lambda_{\mathrm{tor}}, \qquad B_0\bar{\hat\gamma} = -\hat\delta, \qquad B_0\bar{\hat u} = \hat w,
\]
with matrix $\left(\begin{smallmatrix} 0 & 1 \\ -1 & 0\end{smallmatrix}\right)$ in the above identifications; in particular $\det B_0 = 1$, so $B_0$ is unimodular. Moreover $\Lambda^{A_1} := \ker(A_1-I) = \langle \varepsilon,\hat\delta\rangle$ and $\Lambda^{A_2} := \ker(A_2-I) = \langle \varepsilon',\hat\delta\rangle$ are saturated, $\gamma(\varepsilon) = \gamma(\varepsilon') = 1$, and with the twist vectors
\[
v_1 := \varepsilon \in \Lambda^{A_1}, \qquad v_2 := -\varepsilon' \in \Lambda^{A_2}
\]
one has $(\ell_0,\ell_1,\ell_2) := (0,\gamma(v_1),\gamma(v_2)) = (0,1,-1)$.
\end{lemma}

\begin{proof}
By (2.3), for $\lambda = a\hat\gamma+b\hat u+c\hat w+d\hat\delta$,
\[
(M_0-I)\lambda = -N^t\lambda = -a\hat\delta+b\hat w.
\]
Hence $(M_0-I)\lambda = 0$ if and only if $a=b=0$, i.e.\ $\lambda \in \Lambda_{\mathrm{tor}}$; and the image is exactly $\langle \hat w,\hat\delta\rangle = \Lambda_{\mathrm{tor}}$, the pair $(a,b)$ ranging over all of $\dbZ^2$. Thus $M_0-I$ kills $\Lambda_{\mathrm{tor}}$ and the induced map $B_0$ sends $\bar{\hat\gamma}\mapsto -\hat\delta$, $\bar{\hat u}\mapsto \hat w$, i.e.\ $(a,b)^t \mapsto (b,-a)^t$ in the chosen coordinates, of determinant $1$.

For $\lambda$ as above, (2.3) gives
\[
(A_1-I)\lambda = a(6\hat u-6\hat w-2\hat\delta)+b(-\hat u-\hat w+\hat\delta)+c(\hat u-2\hat w),
\]
\[
(A_2-I)\lambda = a(-6\hat w+3\hat\delta)+b(-\hat u+\hat w)+c(-\hat u-\hat w+\hat\delta),
\]
so $(A_1-I)\lambda = 0$ reads $6a-b+c = -6a-b-2c = -2a+b = 0$, i.e.\ $b=2a$, $c=-4a$, $d$ free, and $(A_2-I)\lambda = 0$ reads $b+c = -6a+b-c = 3a+c = 0$, i.e.\ $b=3a$, $c=-3a$, $d$ free. These kernels are saturated, being kernels of endomorphisms of the torsion-free group $\Lambda$, and $\gamma(\varepsilon) = \gamma(\varepsilon') = 1$ by definition of $\gamma$.
\end{proof}

\begin{lemma}
Let $G$ act on $V$ tautologically and on $\Lambda$ through $\mathfrak A$. Then:
\begin{enumerate}[label=(\roman*)]
\item $V^{T_1} = \langle \gamma, 2u+w+3\delta\rangle$, $V^{T_2} = \langle \gamma, u+w+2\delta\rangle$, $V^G = \dbZ\gamma$;
\item $\Lambda^G = \Lambda^{A_1} \cap \Lambda^{A_2} = \dbZ\hat\delta$;
\item every $g \in G$ preserves $\langle \gamma,u,w\rangle$ and acts trivially on $V/\langle\gamma,u,w\rangle$, and every $\mathfrak A(g)$ satisfies $\gamma \circ \mathfrak A(g) = \gamma$; with $\Sigma_V := \sum_{g\in G}\operatorname{im}(g-I)$ and $\Sigma_\Lambda := \sum_{g\in G}\operatorname{im}(\mathfrak A(g)-I)$,
\[
\Sigma_V = \langle \gamma,u,w\rangle, \qquad \Sigma_\Lambda = \ker\gamma,
\]
so the coinvariants $V_G \cong \dbZ$ (generated by the class of $\delta$) and $\Lambda_G \cong \dbZ$ (via $\gamma$) are free of rank $1$; equivalently $[(T_1-I)\mid (T_2-I)]$ and $[(A_1-I)\mid (A_2-I)]$ have Smith normal form $\mathrm{diag}(1,1,1,0)$;
\item the smallest $\langle A_1,A_2\rangle$-invariant subgroup of $\Lambda$ containing $\Lambda_{\mathrm{tor}}$ is $\ker\gamma$.
\end{enumerate}
\end{lemma}

\begin{proof}
(i) For $x = a\gamma+bu+cw+d\delta$, (2.2) gives
\[
T_1 x = (a-6c+2d)\gamma+(-b+c+d)u+(-b+d)w+d\delta,
\]
\[
T_2 x = (a+6b-3d)\gamma+(-c+d)u+bw+d\delta,
\]
so $T_1x = x$ iff $d=3c$, $b=2c$, and $T_2x = x$ iff $d=2b$, $c=b$. Comparing $a\gamma+c(2u+w+3\delta) = a'\gamma+b(u+w+2\delta)$ forces $b=c=0$, so $V^G = \dbZ\gamma$.

(ii) By Lemma~\ref{lem:B0} an element of the intersection is $a\varepsilon+d\hat\delta = a'\varepsilon'+d'\hat\delta$; the $\hat\gamma$- and $\hat u$-coefficients give $a=a'$ and $2a=3a'$, so $a=0$.

(iii) By (2.2) each generator $T_j$ preserves $U := \langle \gamma,u,w\rangle$ and acts trivially on $V/U$; these two properties are stable under composition and inversion, hence hold for every $g \in G$, so $\operatorname{im}(g-I) \subseteq U$ and $\Sigma_V \subseteq U$. Likewise $\gamma \circ A_j = \gamma$ (Lemma 2.4) gives $\gamma \circ \mathfrak A(g) = \gamma$ for every $g \in G$, whence $\operatorname{im}(\mathfrak A(g)-I) \subseteq \ker\gamma$ and $\Sigma_\Lambda \subseteq \ker\gamma$. For the reverse inclusions it suffices to use the generators: $\Sigma_V$ contains the images of the basis vectors under $T_j-I$, namely
\[
-2u-w, \quad -6\gamma+u-w, \quad 2\gamma+u+w, \quad 6\gamma-u+w, \quad -u-w, \quad -3\gamma+u,
\]
and $(-2u-w)-(-u-w) = -u$, then $w$, then $2\gamma$ and $3\gamma$, hence $\gamma$, lie in $\Sigma_V$; while $(A_1-I)\hat w+(A_2-I)\hat u = -\hat w$, then $\hat u$, then $\hat\delta = (A_1-I)\hat u+\hat u+\hat w$ lie in $\Sigma_\Lambda$. Both $\Sigma_V$ and $\Sigma_\Lambda$ are spanned by parts of the bases, hence are direct summands with free rank-$1$ quotients.

(iv) Such a subgroup contains $\hat w$ and $A_1\hat w = \hat u-\hat w$, hence $\hat u$ and $\hat\delta$, i.e.\ $\ker\gamma$; and $\ker\gamma$ is invariant since $\gamma \circ A_j = \gamma$.
\end{proof}

\subsection{The invariant alternating form}

\begin{lemma}\label{lem:Q0}
The group of $G$-invariant alternating forms on $V$ is infinite cyclic, generated by $Q_0$ with
\[
Q_0(\gamma,\delta) = 1, \qquad Q_0(u,w) = 6, \qquad Q_0(\gamma,u) = Q_0(\gamma,w) = Q_0(u,\delta) = Q_0(w,\delta) = 0;
\]
thus $Q_0$ is, up to sign, the unique primitive invariant alternating form. Moreover $b(x,y) := Q_0(x,Ny)$ is symmetric and descends to $V/\ker N \cong \dbZ^2$; in the basis given by the classes of $w,\delta$ its Gram matrix is
\[
\mathrm{diag}(6,-1),
\]
so $b$ is nondegenerate of signature $(1,1)$ and discriminant $-6$.
\end{lemma}

\begin{proof}
Write $q_{xy} := Q(x,y)$ for an invariant alternating $Q$. By (2.2), $T_1$-invariance on the pairs $(\gamma,u)$ and $(\gamma,w)$ reads
\[
-q_{\gamma u}-q_{\gamma w} = q_{\gamma u}, \qquad q_{\gamma u} = q_{\gamma w},
\]
i.e.\ $2q_{\gamma u}+q_{\gamma w} = 0$ and $q_{\gamma u} = q_{\gamma w}$, whence $3q_{\gamma u} = 0$ and
\[
q_{\gamma u} = q_{\gamma w} = 0.
\]
Using these first two equations to drop all terms in $q_{\gamma u}, q_{\gamma w}$ from the remaining expansions, $T_1$-invariance on the pairs $(u,\delta)$ and $(w,\delta)$ reads
\[
-q_{u\delta}-q_{w\delta} = q_{u\delta}, \qquad -6q_{\gamma\delta}+q_{uw}+q_{u\delta} = q_{w\delta},
\]
and, in the same way (again using the first two equations), $T_2$-invariance on the same four pairs reads
\[
q_{\gamma w} = q_{\gamma u}, \qquad -q_{\gamma u} = q_{\gamma w}, \qquad -q_{u\delta} = q_{w\delta}, \qquad 6q_{\gamma\delta}-q_{uw}+q_{w\delta} = q_{u\delta};
\]
the remaining pairs give no condition. Hence $q_{\gamma u} = q_{\gamma w} = q_{u\delta} = q_{w\delta} = 0$ and $q_{uw} = 6q_{\gamma\delta}$, i.e.\ $Q \in \dbZ Q_0$, and $Q_0$ is invariant. For $b$: $\ker N = \operatorname{im} N = \langle \gamma,u\rangle$ and $Q_0$ vanishes on $\langle\gamma,u\rangle \times \langle\gamma,u\rangle$, so $b(x,y) = 0$ whenever $x$ or $y$ lies in $\ker N$, and $b$ descends; finally
\[
b(w,w) = Q_0(w,-u) = 6, \qquad b(\delta,\delta) = Q_0(\delta,\gamma) = -1, \qquad b(w,\delta) = Q_0(w,\gamma) = 0 = Q_0(\delta,-u) = b(\delta,w). \qedhere
\]
\end{proof}

\begin{remark}
In \S4 it is the unimodularity of $B_0$ that is used: the lattice $B_0\bar\Lambda$ of cocharacter translations is all of $\Lambda_{\mathrm{tor}}$, and this is what makes the toric filling have an irreducible central fibre $W$ containing a single copy of $\mathrm{dP}_6$; for $|\det B_0| = k > 1$ the same construction would produce $k$ such pieces. The form $Q_0$ is used again in Remark 3.23, in \S7 through the invariant bivector $\eta = u\wedge w+6\,\gamma\wedge\delta$ dual to $Q_0$, and in Appendix A. It is recorded here already in order to place the present data in context: for a one-parameter degeneration of polarized abelian varieties with unipotent monodromy $\exp N$ and polarization $Q$ the form $x \mapsto Q(x,Nx)$ on $V/\ker N$ is definite (\cite{Cle77}, \cite{Mum72}), whereas by Lemma~\ref{lem:Q0} the pair $(Q_0,N)$ has signature $(1,1)$.
\end{remark}

\subsection{The base orbifold, the group $\Delta$, and the representation $\rho_V$}

With the conventions of page 1 (Setup): $B = \dbP^1$ with coordinate $t$, $p_1\colon t=0$, $p_2\colon t=1$, $p_0\colon t=\infty$ with $t_c = 1/t$, $B^\circ = B \setminus \{p_0,p_1,p_2\}$, $m_1=3$, $m_2=4$, $\zeta_j = e^{-2\pi i/m_j}$. Let $\mathcal O$ be the orbifold with underlying space $B$, cone points of orders $3$ at $p_1$ and $4$ at $p_2$, and a puncture at $p_0$, i.e.\ of signature $(0;3,4,\infty)$.

\begin{lemma}
Let $g \in \mathrm{PSL}_2(\dbR)$ be elliptic of order $m < \infty$ with fixed point $z_0 \in \hb$, and $c(z) := (z-z_0)/(z-\bar z_0)$. Then $c$ maps $\hb$ biholomorphically onto the unit disc with $c(z_0) = 0$, and $c \circ g \circ c^{-1}$ is the rotation $s \mapsto \zeta s$ with $\zeta = g'(z_0)$ a primitive $m$-th root of unity; in particular $c \circ g = \zeta \cdot c$.
\end{lemma}

\begin{proof}
$c$ is a M\"obius map carrying $\dbR \cup \{\infty\}$ to the unit circle and $z_0$ to $0$; the disc automorphism $c\circ g \circ c^{-1}$ fixes $0$, hence is $s \mapsto \zeta s$ with $\zeta$ its derivative at $0$, equal to $g'(z_0)$, of order $m$.
\end{proof}

\begin{proposition}\label{prop:uniformisation}
There exist a copy $\hb_z$ of the upper half-plane, a Fuchsian group $\Delta \subset \mathrm{PSL}_2(\dbR)$ and a holomorphic surjection $\pi \colon \hb_z \to B \setminus \{p_0\}$ such that:
\begin{enumerate}[label=(\roman*)]
\item $\Delta$ acts properly discontinuously, $\pi$ induces $\hb_z/\Delta \cong B\setminus\{p_0\}$, $\pi$ is a local biholomorphism off $\pi^{-1}\{p_1,p_2\}$ and has local degree $m_j$ at each point of $\pi^{-1}(p_j)$; the stabiliser of such a point is cyclic of order $m_j$ and all other stabilisers are trivial;
\item there are $z_j \in \pi^{-1}(p_j)$ and $g_1,g_2 \in \Delta$ with $g_j(z_j) = z_j$, $g_j$ of order $m_j$ and $g_j'(z_j) = \zeta_j$, such that
\[
\Delta = \langle g_1,g_2 \mid g_1^3 = g_2^4 = 1\rangle \cong \dbZ/3 * \dbZ/4,
\]
and such that $g_0 := (g_1g_2)^{-1}$ is parabolic; by construction $g_1g_2g_0 = 1$;
\item if $* \in \partial\hb_z$ is the fixed point of $g_0$, its stabiliser in $\Delta$ is exactly $\langle g_0\rangle$, and for all small $r>0$ the component $U_r$ of $\pi^{-1}(\{0<|t_c|<r\})$ whose closure contains $*$ is a $\langle g_0\rangle$-invariant, connected and simply connected cusp neighbourhood of $*$ squeezed between two horodiscs at $*$, i.e.\ $H'' \subseteq U_r \subseteq H'$ for suitable horodiscs $H'', H'$ at $*$; shrinking $r$ so that $U_r$ lies in a strictly smaller horodisc at $*$, the set $U_r$ is closed in $\pi^{-1}(\{0<|t_c|<r\})$ and hence is a connected component of it, every other component lying in a translate $gH'$ with $g \in \Delta \setminus \langle g_0\rangle$; it satisfies $gU_r \cap U_r = \emptyset$ for $g \in \Delta \setminus \langle g_0\rangle$, and $\pi$ induces a biholomorphism $U_r/\langle g_0\rangle \cong \{0<|t_c|<r\}$.
\end{enumerate}
\end{proposition}

\begin{proof}
Since $\tfrac13+\tfrac14 < 1$, $\mathcal O$ is hyperbolic, hence uniformised by a Fuchsian group of signature $(0;3,4,\infty)$ (\cite[Ch.\ 13]{Thu97}, \cite[Ch.\ 10]{Bea83}); explicitly, $\Delta$ is the orientation-preserving subgroup of the group generated by the reflections $\sigma_a,\sigma_b,\sigma_c$ in the sides $a=BC$, $b=CA$, $c=AB$ of the hyperbolic triangle $T$ with vertices $A,B,C$ counterclockwise and angles $\pi/3,\pi/4,0$ (Poincar\'e's theorem; \cite[Ch.\ 7, 9--10]{Bea83}). The quotient surface is a once-punctured sphere, i.e.\ $\dbC$; composing with the affine automorphism carrying the order-$3$ cone point to $0=p_1$ and the order-$4$ one to $1=p_2$ yields $\pi$, and (i) is the standard description of a Fuchsian group by its signature (\cite[Ch.\ 10]{Bea83}).

(ii) Put $x := \sigma_b\sigma_c$, $y := \sigma_c\sigma_a$, $z := \sigma_a\sigma_b$, so $xyz=1$ and $\Delta = \langle x,y,z \mid x^3=y^4=1, xyz=1\rangle$ (\cite[Ch.\ 10]{Bea83}). The relevant rule is: if two geodesics $\alpha,\beta$ meet at $P$ and $\theta$ denotes the angle from $\alpha$ to $\beta$ measured counterclockwise, then the product $\sigma_\beta\sigma_\alpha$ (reflect in $\alpha$ first, then in $\beta$) is the counterclockwise rotation about $P$ through $2\theta$ (\cite[Ch.\ 7]{Bea83}). Since the vertices $A,B,C$ are counterclockwise, the counterclockwise angle at $A$ from $c=AB$ to $b=CA$ is the interior angle $\pi/3$, and that at $B$ from $a=BC$ to $c=AB$ is $\pi/4$; hence $x=\sigma_b\sigma_c$ is the counterclockwise rotation about $A$ through $2\pi/3$, $y=\sigma_c\sigma_a$ the counterclockwise rotation about $B$ through $2\pi/4$, and $z$ is parabolic, fixing the ideal vertex $C$; by the previous lemma, $x'(A) = e^{2\pi i/3}$ and $y'(B) = e^{2\pi i/4}$. With
\[
z_1 := A, \qquad z_2 := B, \qquad g_1 := x^{-1}, \qquad g_2 := y^{-1}
\]
the two rotations are reversed simultaneously, so that $g_1'(z_1) = \zeta_1$ and $g_2'(z_2) = \zeta_2$, of orders $3,4$. Inverting $xyz=1$ gives $g_2g_1 = z$, so $g_1g_2 = g_1zg_1^{-1}$ is parabolic, hence so is $g_0$; eliminating $z$ gives the stated presentation.

(iii) The stabiliser of $*$ is $\langle g_0\rangle$. By the standard description of a cusp (\cite[Ch.\ 10]{Bea83}) the stabiliser $\Delta_*$ of the fixed point of a parabolic element of a Fuchsian group is infinite cyclic, generated by a primitive parabolic $h$, and there is a horodisc $H'$ at $*$ with $gH' \cap H' = \emptyset$ for $g \notin \Delta_*$. Write $g_0 = h^n$ with $n \ge 1$; we claim $n=1$, i.e.\ that $g_0$, equivalently $g_0 = g_1^{-1}g_2^{-1}$, is not a proper power in $\Delta$. Indeed, in the free product $\Delta = \langle g_1\rangle * \langle g_2\rangle$ the element $g_1g_2$ is cyclically reduced of length $2$, and cyclically reduced length is a conjugacy invariant. An element $h \in \Delta$ of finite order is conjugate into a free factor, so all its powers have finite order and cannot be parabolic; hence $h$ has infinite order and is conjugate to a cyclically reduced word $h_0$ of even length $\ell \ge 2$, so that $h^n$ is conjugate to the cyclically reduced word $h_0^n$ of length $n\ell$. If $n \ge 2$ this length is $\ge 4 > 2$, a contradiction. Thus $n=1$ and $\Delta_* = \langle g_0\rangle$, and $H'/\langle g_0\rangle$ embeds in $\hb_z/\Delta$ as a punctured-disc neighbourhood of the puncture.

\emph{The cusp coordinate.} Conjugate in $\mathrm{PSL}_2(\dbR)$ so that $* = \infty$ and $H' = \{\operatorname{Im}\zeta > C\}$; then $g_0(\zeta) = \zeta+\epsilon_0$ with $\epsilon_0 \in \dbR^\times$, and after the further conjugation $\zeta \mapsto \zeta/|\epsilon_0|$ we may assume $\epsilon_0 \in \{+1,-1\}$. The sign $\epsilon_0$ is intrinsic, the elements of $\mathrm{PSL}_2(\dbR)$ fixing $\infty$ being the maps $\zeta \mapsto a\zeta+c$ with $a>0$, whose conjugation action preserves it. Its value plays no role in the present proof, since $q := e^{2\pi i\zeta}$ is $\langle g_0\rangle$-invariant for either sign and identifies $H'/\langle g_0\rangle$ with a punctured disc $D^*$; the value is determined once the period map $\tau$ is taken into account, and Remark~2.12 gives $\epsilon_0 = -1$. Thus $\pi$ induces an injective holomorphic map $\psi\colon D \to B$ with $\psi(0) = p_0$, where we read $B$ near $p_0$ in the coordinate $t_c$ and set $\psi(0)=0$. Choose $r_0>0$ with $\{|t_c|<r_0\} \subseteq \psi(D)$, and let $0<r<r_0$. Then $V_r := \psi^{-1}(\{|t_c|<r\})$ is a simply connected neighbourhood of $0$ in $D$, biholomorphic to a disc, and
\[
U_r = \{\zeta \in H' : e^{2\pi i\zeta} \in V_r \setminus\{0\}\}
\]
is the component of $\pi^{-1}(\{0<|t_c|<r\})$ accumulating at $*$: it is the preimage under $\zeta \mapsto e^{2\pi i\zeta}$ of a punctured simply connected domain around $0$, hence connected, simply connected and $\langle g_0\rangle$-invariant, and $\pi$ induces $U_r/\langle g_0\rangle \cong \{0<|t_c|<r\}$; disjointness of its $\Delta$-translates off $\langle g_0\rangle$ follows from that of $H' \supseteq U_r$. It is moreover the only component accumulating at $*$: since $\pi$ is the quotient by $\Delta$ and $U_r$ is a component of $\pi^{-1}(\{0<|t_c|<r\})$ with trivial setwise stabiliser outside $\langle g_0\rangle$, every component of $\pi^{-1}(\{0<|t_c|<r\})$ is a translate $gU_r$ with $g \in \Delta$, and every component other than $U_r$ is such a translate with $g \notin \langle g_0\rangle$, hence lies in $gH'$; since $U_r \subseteq H'$ is contained in a horodisc at $*$, the closure of $gU_r$ meets $\partial\hb_z$ only at $g*$, and $g* \ne *$ because the stabiliser of $*$ is exactly $\langle g_0\rangle$. Hence no other component accumulates at $*$. Finally, choosing $0<\varrho<\varrho'$ with $\{|q|<\varrho\} \subseteq V_r \subseteq \{|q|<\varrho'\}$ gives horodiscs $H'' = \{\operatorname{Im}\zeta > \log(1/\varrho)/2\pi\} \subseteq U_r$ and $U_r \subseteq \{\operatorname{Im}\zeta > \log(1/\varrho')/2\pi\} \cap H'$, as claimed.
\end{proof}

\begin{remark}[The translation sign at the cusp]
The normalisation $g_0(\zeta) = \zeta+\epsilon_0$, $\epsilon_0 \in \{\pm1\}$, of the proof of Proposition~\ref{prop:uniformisation}(iii) concerns the uniformising coordinate $\zeta$ on $\hb_z$, whereas the relation $\tau \circ g_0 = \tau-1$ of \S3 concerns the period map $\tau\colon \hb_z \to \hb$ and expresses its equivariance along $\rho_\Lambda$, i.e.\ $\rho_\Lambda(g_0) = M_0$. The two maps are distinct, but the second relation determines the first sign: we show that $\epsilon_0 = -1$.

Indeed, $\tau$ is equivariant along the surjection $\Delta \to \mathrm{PSL}_2(\dbZ)$ of Remark~2.13, so $j\circ\tau = 1728\,t\circ\pi$ on $\hb_z$; equivalently, on $U_r$ the function $e^{2\pi i\tau}$ is, through the identifications of Proposition~\ref{prop:uniformisation}(iii), the local coordinate $q$ up to a unit \emph{(rest of argument continues on p.\,13 of the source, not yet transcribed)}.
\end{remark}
\newpage


\begin{center}
\emph{[Remainder of the document (\S\S3--10 and Appendices A--B, pp.\ 13--108) not yet transcribed.]}
\end{center}

\end{document}
